\PBchapter{Estudos}

%\PByear{1997}

Em 1997, resolvi que precisava voltar a estudar. 

Faço cursinho.

Passei em uma universidade pública estadual no curso de Matemática, em quadragésimo lugar de quarenta vagas.

Concluo o curso? Não.

Em 2001, a universidade entra em greve. 

Greve essa que dura “apenas” 169 dias. 

Aqui eu preciso, não tem jeito, fazer um pequeno relato. 

Desde que entrei nessa universidade, tivemos greve em 1998, aproximadamente 30 dias, 1999, aproximadamente 40 dias e 2000, aproximadamente 60 dias. 

Ou seja, eu já tinha percebido uma tendência, já que estava envolvido, com Centro Acadêmico de Matemática e Diretório Central de Estudantes.

Porém, na assembleia de 2001, que decidiu pelo apoio à greve, eu fui contra, pois já tinha esse histórico recente. 

E era um período de muita instabilidade, e tinha certeza que uma greve comum, não resolveria o problema e acabaria sendo muito mais extensa. 

Aqui é um ponto que pela primeira vez eu pensei em mim, pois a empresa que eu trabalhava, a única exigência era estar estudando.

Então chegamos a 2002, e um ponto de virada. Eu afirmando nas reflexões da adolescência, que não tinha pontos de virada.

Como relatei, eu precisava estudar. 

Eu não queria fazer cursinho ou algo assim. Isso resolveria meu problema na empresa. 

Mas na minha cabeça eu estava regredindo, dando pessoas para trás - bom, aqui vou confessar, não era impressão, eu quase sempre estava dando passos para trás como vão ver mais a frente. 

Então resolvi fazer vestibular para uma faculdade particular. Matemática ou curso correlato?

Claro que não. 

Resolvi prestar para Direito. 

Passei? Também não. 

Mas, por ser uma faculdade particular, me chamaram para fazer Processamento de Dados. 

Na empresa em que eu atuava, já estava meio que inserido nesse meio, então resolvi fazer.

Concluo o curso? 

Gente… caros… vocês já sabem a resposta! 

Claro que não. 

Em 2002, fui demitido. 

Em 2003 minha mãe ficou desempregada.  

Em 2003 consigo um estágio em uma empresa de software, e dá um alívio para continuar no curso. 

Mas em 2004 fui demitido.

Até consigo um trabalho logo em seguida, mas as contas só aumentam e não consigo me dedicar ao curso. 

No final de 2004 tinha TCC para fazer, e acabei fazendo TCC de outras pessoas, não só do curso, mas para pessoas de outras universidades e cursos, para ganhar dinheiro. 

Na graduação fico devendo apenas o trabalho de conclusão de curso. 

Até tento me matricular em 2005 para concluir o TCC, mas devido aos débitos financeiros, não consigo.

Escrevendo e lendo esse meu relato, a sensação é de desespero.

Uma vontade de voltar e dar um belo pescotapa em mim mesmo.

Porque vocês imaginam que a saga de estudo termina aqui?

Logíco que não.

De 2005 até meados de 2012, é só trabalho. 

Sem crescimento profissional e financeiro. 

Só vivendo. 

Tomando algumas decisões pouco pensadas, como por exemplo, arrendar um bar, e trabalhar depois do trabalho. Lógico que não deu certo.

Então no ano de 2012, resolvi que precisava voltar a estudar.

A empresa que eu estava, era uma revenda de um software acadêmico. 

Então, fui e fiz uma redação para vagas remanescentes em uma dessas faculdades que a empresa atendia. 

Entrei no curso de Tecnólogo em Análise e Desenvolvimentos de Sistemas, vulgo TADS. 

Seria um filho do curso de Processamento de Dados? 

Consigo toda a documentação desse curso de Processamento de Dados e consigo eliminar um monte de disciplinas, indo para o quinto período do curso. 

Como o curso era de oito períodos, então, na teoria, mais três períodos, ou seja, um ano e meio, eu concluiria. 

Pergunto novamente, concluo o curso? 

Vocês já sabem que não.

Na segunda semana de aula, eu exausto do dia de trabalho, chego em uma aula de comércio eletrônico, o professor pede pra baixar um livro que foi tirado fotos das páginas para a gente ler. Sim, fotos.

Ele chegou a olhar pra mim, mas não teve coragem de pedir para eu ler. 

Então ele vendo que estava esquisito, para a atividade, pede para gente pesquisar sobre comércio eletrônico. Aqui precisam entender que o ano é início de 2012 e estou com 35 anos. 

A média da sala não passa de 23 anos. E sim, eu era mais velho que a maioria dos professores.

Nesse momento, eu pesquiso: “Curso superior de computação”. 

Vou tentar resumir muito, mas já perceberam que o resumo não é bem o meu forte. 

Eu encontrei o curso com o mesmo nome em uma Universidade Federal, que tinha um campus em Medianeira, cidade do Paraná, cerca de 75 km de Cascavel. 

E encontro um edital de vagas remanescentes para o curso e as inscrições se encerram na sexta da mesma semana, e só para entenderem, isso foi numa quarta feira. 

Detalhe, apenas três vagas. 

Nem penso muito, saio da sala, vou na secretária, pois precisava daquela documentação que tinha levado duas semanas antes. 

Como não ia dar tempo de enviar pelos correios, agendei uma reunião para sexta-feira em uma faculdade particular que a empresa que eu trabalhava, atendia em Medianeira. 

Levo toda a documentação na Universidade. 

Assim, entendam, eu não sabia nem da existência dessa Universidade. 

Aqui podem imaginar como foi o meu final de semana. 

O resultado iria sair somente na quinta-feira da outra semana. 

Pois bem, sexta-feira de manhã, a primeira coisa foi acessar o site da universidade.
 
Certo, sem suspense, eu passei. Em terceiro, empatado com o quarto colocado. O desempate foi a idade. 

Então começo o segundo semestre de 2012 em uma universidade federal. 

E como nem tudo na vida são flores, um mês depois, o que acontece? 

Greve. 

Eu já tinha montado um esquema para esse primeiro semestre do curso. 

Uma fase de adaptação, onde concentrei minhas disciplinas segunda e terça, então eu ia no domingo, e ficava em uma pensão, e retornava na quarta feira de manhã para Cascavel. 

Nesses dias, segunda e terça, eu trabalhava remotamente. 

A greve durou 32 dias, assim, tivemos aulas em janeiro e fevereiro. 

Vem o ano de 2013. 

Vem um certo descontentamento com a vida. 

E o que a gente faz quando está descontente e sem perspectiva? 

Para? 

Avalia a situação? 

Claro que não. 

O negócio é montar um hostel em Foz do Iguaçu com mais dois amigos. 

Mas que excelente ideia não. 

Durou de maio de 2013 a julho de 2014. 

Foram os melhores e piores meses da minha vida. 

Tranquei a faculdade, me demiti do meu trabalho. 

Tive experiências absurdas. 

Recebi pessoas de aproximadamente 42 países. 

Troca de cultura. 

Conheci o mundo, sem sair do lugar. 

Isso teve um custo. 

Mas eu ainda não consigo me arrepender. 

Talvez, não se arrepender seja um problema. Mas simplesmente não consigo.

Eu acredito que não seria a pessoa que sou hoje, se não tivesse tido essa experiência.

Sei que ainda parece que estou enrolando. 

Então vamos seguindo. 

Volto para Cascavel, desempregado, com quase 37 anos, sem perspectiva. 

Arrumo um trabalho. 

Mas o descontentamento é grande. 

Resolvi então, reabrir o curso na universidade federal em Medianeira. 

Consegui na bacia das almas. Explico. 

O curso de TADS, estava sendo extinto, e isso já era de meu conhecimento em 2012. 

Logo eu entro no curso de TADS para garantir a vaga, mas no semestre seguinte eu precisava fazer a mudança de curso. 

Me mudo pra Medianeira novamente. 

Retorno para a mesma república que tinha ficado de 2012 a 2013, somente com o dinheiro de um mês, e um pouco para comida. 

E saio de porta em porta em Medianeira em busca de trabalho. 

E pasmem, em 2 dias eu consegui. 

Empresa de desenvolvimento de software para lojas de pequeno porte.

Como eu disse, eu precisava mudar de curso, e isso tem um prazo, tem edital. 

Universidade federal o que vale é o que está no papel. 

E eu anoto bem ‘erradinho’ dia 13/11/2014. 

Mas a data correta era dia 12 de novembro às 13 horas. 

Chegou então o dia 13 de novembro, vou na secretaria para realizar a mudança de curso, e só ouço a secretária dizendo: “O prazo era até ontem”.

Aqui é daqueles momentos da vida, que você não tem reação. 

Na minha mente veio a imagem de uma rota de ônibus que sempre via na rodoviária quando ia de Medianeira para Cascavel e vice-versa.
 
Essa rota saia de Cruz Alta, no Rio Grande do Sul, e o destino era Barreiras, na Bahia. 

Cheguei a abrir a carteira, e contar, R\$ 167,00, era o que eu tinha pra passar os últimos 25 dias até receber o próximo salário. 

Pensei, até que trecho isso daria dessa rota? 

Devo ter ficado parado por uns 15 a 20 segundos. 

Até que percebo a secretária falando comigo, e pedindo pra eu enviar um e-mail para o contato que ela tinha anotado em um papel. 

Sim, eu consegui efetuar a transferência de curso e garantir minha vaga. Mas precisava relatar esse momento.

Como prometi, aqui vou resumir. Ou não. Me formei em Bacharel em Ciência da Computação em agosto de 2019. 

Eu menti. Não consigo resumir.  

Teve um monte de alegrias, conquistas e contratempos de 2015 a 2019.

Preciso relatar um deles, para uma das lições mais sensacionais que recebi na vida.

Coloquei na cabeça que iria fazer dupla diplomação, pois o Curso/Universidade tinha uma parceria com uma instituição de Portugal. 

Faço minha inscrição, e até então, eram eu e mais um inscrito. Como eram, acho, seis vagas, então, estava tranquilo. 

No último dia de prazo, oito alunos se inscrevem.

E lógico, minha pontuação não foi o suficiente para competir com quem teve uma trajetória regular no curso.

Estou no último período. Às vésperas de apresentar o TCC.

Sobre o TCC, existia uma exigência no curso de enviar a publicação de um artigo para algum evento. 

Envie. O artigo foi aceito. E fui convidado para apresentá-lo em formato painel no evento.

Dias depois, o artigo caiu pra cima. O painel passou a ser uma apresentação. Auditório.

O evento foi em SP. Mais uma saga. Eu não tinha grana pra ir.

Vaquinha virtual.

E vem um apoio inesperado, um ex-aluno da UTFPR, e colega que eu tinha estudado, estava em SP trabalhando em uma empresa de consultoria.

Ele consegue o patrocínio, que não ajuda só com a minha viagem, mas de mais 2 colegas. 

Apresentei o TCC no evento e nessa empresa.

Como eu disse, estava no final do curso, o TCC faltava ser concluído.

Então, chego para um professor, que inclusive foi nesse evento em SP, e comentei que vou segurar o TCC pra tentar novamente o intercâmbio em Portugal. Ele sabia um pouco da minha trajetória de estudo.

É então que ele calmamente diz:

- Eu conheço um pouco da sua trajetória. Você precisa fechar ao menos um ciclo, concluir algo. Depois se quiser iniciar um  novo ciclo, sem problema. Mas está com 42 anos. Precisa terminar o que começa.

Cirúrgico.

Minha reação foi não ter reação. 

Por uns cinco minutos.

Então me internei em um laboratório. 

Escrevi meu TCC.

Nesse dia, em casa, fui até às quatro da manhã escrevendo.

Depois virei mais uma noite. 
E conclui.

Não passei com honras. 

Não foi o melhor TCC. 

Mas, conclui um ciclo.